% !TeX root = xpsejal00-gcc-plugin.tex
\documentclass[xpsejal00-gcc-plugin.tex]{subfiles}
\begin{document}

\chapter{Introduction}\label{chap:introduction}

This bachelor's thesis focuses on the design and implementation of a modernized reimplementation of the Code Listener infrastructure, a C++ library that serves as a vital bridge between compilers and static analysis tools, particularly the Predator shape analysis tool suite.
The original Code Listener architecture, while successful in enabling static analysis of C programs that manipulate heap memory, exhibits several design limitations that have accumulated over time.
This thesis aims to address these constraints by replacing the callback-driven pipeline with a persistent Code Model, demoting analyses from hard-wired pre-analysis steps to optional, on-demand annotation services, and introducing a structured query interface and uniform support for exporting results in common formats.

\section{Motivation}\label{sec:motivation}

Software development is a complex and error-prone process, especially as software systems grow in size and complexity.
As software systems evolve, certain categories of defects become undetectable solely through testing, as test suites can only evaluate explicitly covered execution paths, leaving much of the program's state space unexplored.

Nowhere is this limitation more critical than in low-level system programming languages such as C and C++, where developers bear the burden of manual memory management.
In these environments, insidious bugs such as memory leaks, buffer overflows, use-after-free errors, and dangling pointer dereferences frequently hide within unexplored execution paths.
The consequences of these hidden defects can be catastrophic, particularly in mission-critical domains like aerospace engineering.
While rigorous testing prevents many fatal accidents, memory mismanagement has repeatedly compromised modern aircraft.
For instance, regulatory agencies have had to issue mandatory airworthiness directives for both the Boeing 787 and the Airbus A350, requiring airlines to periodically reboot the aircraft to clear memory buffers and prevent the total mid-flight loss of critical systems.
These near-misses underscore a fundamental reality: traditional dynamic testing is insufficient to guarantee memory safety in high-stakes environments.

To prevent such failures, developers must rely on advanced static analysis tools, such as the Predator shape analysis suite, which is capable of formally verifying a program's state space without executing it.
However, developing and maintaining these analyzers presents a significant architectural challenge: extracting the program's semantics from the compiler's internal representation.
Analyzers must interface with compilers to understand the code, but native compiler APIs are notoriously complex, tightly coupled to specific versions, and largely undocumented.
Building analysis tools directly on top of these raw APIs leads to brittle, maintenance-heavy software.

To overcome this bottleneck, a clean abstraction layer is required.
Specifically, researchers need a modular middleware that can abstract away the massive, often undocumented codebases of industrial compilers, providing instead a concise, unified, and well-documented object-oriented API.\@
This ensures that static analysis tools can be developed independently of the underlying compiler architecture.

% TODO: some of this may be in the abstract:
% Developing correct software presents significant challenges.
% As software systems grow in size and complexity, certain categories of defects become undetectable solely through testing.
% Test suites can only evaluate explicitly covered execution paths, leaving much of the program's state space unexplored.
% Formal verification methods, particularly static program analysis, address this limitation by reasoning about program behavior without execution, thereby examining all reachable states or proving properties that hold unconditionally.
%
% C programs that manipulate heap memory using low-level pointer operations, such as dynamic linked lists, trees, and graph-like structures with arbitrary run-time shapes, present significant challenges for static analysis.
% Errors in such code, including dangling pointer dereferences, double frees, and memory leaks, are difficult to detect and can have severe consequences, particularly in safety-critical domains such as aerospace engineering.
% The Predator shape analyzer~\cite{PredatorSuite2016} and the related Forester tool~\cite{CodeListener2012} were specifically developed to detect this class of defects.
%
% Both tools are built on top of the \emph{Code Listener} (CL) infrastructure~\cite{CodeListener2012}, a C++ library that acts as a bridge between an industrial-strength compiler and the analyzers.
% By implementing Predator and Forester as \emph{GCC plug-ins}~\cite{GCCPluginAPI}, the developers gained a crucial property: whatever source file GCC can compile, the analyzer can also process, meaning that there is no risk that the static analysis tool will fail to parse code that the production compiler accepts.
% Moreover, Code Listener's interface to analyzers is completely independent of GCC, allowing the same analyzer to be attached to a different compiler without changing the analysis code.
% This independence was later demonstrated by the creation of an LLVM/Clang adapter~\cite{LLVMAdapter2014} and a Sparse adapter~\cite{SparseAdapter2012} in two bachelor's theses.
%
% Despite notable achievements, including Predator's gold medals at the international software verification competition SV-COMP and its application to industrial code~\cite{CodeListener2012}, the original Code Listener architecture exhibits several design limitations that have accumulated over time.
% Code is translated from GCC's GIMPLE intermediate representation into a stream of \emph{listener callbacks}, which pass through a chain of filters—most prominently, the \emph{switch-to-if} transformation that rewrites all \texttt{SWITCH} instructions into conditional branches—before being accumulated by the \emph{code storage} component.
% Prior to analyzer inspection, four analyses are executed unconditionally: call graph construction, loop-closing edge annotation, points-to analysis, and liveness analysis.
% These transformations and analyses are opaque to analyzers and cannot be disabled, even when unnecessary.
% Furthermore, the architecture lacks standardized mechanisms for analyzers to attach custom annotations to the code model, a structured query interface, and uniform support for exporting results in common formats, such as JSON.
%
% The primary focus of this thesis is \emph{research and engineering}, specifically the design and implementation of a modernized reimplementation of Code Listener that addresses the aforementioned constraints.
% The new implementation, referred to throughout this text as the \emph{New Code Listener}, replaces the callback-driven pipeline with a persistent \emph{Code Model} that faithfully captures the GCC GIMPLE IR without requiring mandatory transformations.
% Analyses are demoted from hard-wired pre-analysis steps to optional, on-demand \emph{annotation services} that the analyzer requests explicitly via a type-safe \emph{Analysis Manager}.
% A~\emph{Query System} exposes the Code Model through a C++ range-based API.
% A~\emph{Diagnostic Reporter} lets analyzers emit warnings and errors in GCC's standard format.
% Optional \emph{Exporters} can serialize the Code Model to JSON, a Graphviz visualization, or a human-readable three-address code listing.
% Two standalone command-line tools, \texttt{cl\_merge} and \texttt{cl\_analyze}, enable whole-program analysis without re-invoking GCC.
%
% Backward compatibility with Predator is maintained through a dedicated \emph{Predator Compatibility Layer}, a thin adapter that translates the new API into the data layout expected by the existing Predator kernel, allowing the regression test suite to run against the new infrastructure.

\section{Thesis structure}\label{sec:thesis_structure}

The remainder of this thesis is structured as follows.
% Chapter~\ref{chap:sota} provides a detailed description of the existing Code Listener infrastructure, specifically its intermediate representation, the callback-based API, the code storage model, and the analyses performed before control is handed to an analyzer.
Chapter~\ref{chap:sota} provides a deeper technical analysis of the legacy Code Listener architecture.
% Chapter~\ref{chap:design} presents the architecture and design of the New Code Listener, covering the object-oriented class hierarchy, the Adapter pattern used for compiler integration, and the interface to GCC.
Chapter~\ref{chap:design} outlines the architectural modernization, detailing the transition from a callback-driven pipeline to a persistent Code Model.
% Chapter~\ref{chap:implementation} documents the development environment, portability considerations, and the resulting code metrics.
Chapter~\ref{chap:implementation} documents the implementation process, including development environment setup, portability considerations, and code metrics regarding the refactored infrastructure.
% Chapter~\ref{chap:testing} evaluates the implementation through integration tests and a regression run against Predator's test suite.
Chapter~\ref{chap:testing} presents the methodology and results of regression testing, demonstrating the correctness and performance of the new implementation against the existing test suite. % TODO: along with integration tests that validate the new architecture's features.
% Chapter~\ref{chap:conclusion} summarizes the results and outlines directions for future work.
Finally, Chapter~\ref{chap:conclusion} summarizes the contributions of this work and outlines potential avenues for future research and development.

\end{document}
